%% toc: P15 166 166.6
\chapter{Maximum-Entropy Inference under Declared Constraints}
\label{sec:p15-w02b-p15-043-maxent}

Maximum-entropy inference assigns a probability distribution when the
available information fixes some expectations but does not fix the
distribution. Its inputs are a declared space of possibilities, a declared
reference measure on that space, declared statistics and declared values
for their expectations. Its output is the distribution of largest entropy
among those compatible with the declarations, when such a distribution
exists. This chapter states that optimization problem in general, derives
its solution together with the conditions under which a solution exists,
is unique and has multipliers, and then reconstructs the normative argument
that Jaynes gave for using the solution as an inference
\cite[pp.~620--623, 630]{Jaynes1957InformationTheory}.

Two kinds of claim are kept apart throughout. The optimization results are
theorems about a declared problem: they hold whatever the declarations
mean. The inferential reading, that the optimizer is the distribution one
\emph{ought} to adopt, is a further claim that depends on premises about
probability, uncertainty and information. A change of reference measure or
of constraint changes the problem and, in general, its answer; such a change
is a change of premises, not new evidence.

The bounded single-mean case appears as
\cref{thm:p15-measurement-asymmetric-model} and the unconstrained finite
case as \cref{thm:p15-rat41-indifference}. Their entropy-maximization
statements are special cases of the results below. A countable form of the Gibbs inequality used here is
\cref{thm:p15-rmdl-gibbs}; the measure version needed for a general
reference measure is proved again in the annex.

\section{The declared problem}
\label{sec:p15-w02b-p15-043-problem}

\begin{hbdefinition}{def:p15-w02b-p15-043-problem}{A declared maximum-entropy problem}
A \emph{declared maximum-entropy problem} consists of the following data.
\begin{enumerate}
\item A \emph{sample space with reference measure}
\((\mathcal X,\mathcal A,\nu)\): a measurable space with a
\(\sigma\)-finite measure \(\nu\) satisfying \(\nu(\mathcal X)>0\).
In the \emph{finite case} \(\mathcal X\) is finite,
\(\mathcal A\) is its power set and \(\nu\{x\}=w(x)>0\) for every
\(x\); counting measure is the case \(w\equiv1\).
\item \emph{Constraint functions} \(f=(f_1,\ldots,f_m)\), each a
measurable map \(f_k:\mathcal X\to\mathbb R\), with \(m\geq0\).
\item \emph{Declared expectation values} \(c=(c_1,\ldots,c_m)\in\mathbb R^m\).
\end{enumerate}
A probability measure \(P\) on \((\mathcal X,\mathcal A)\) is
\emph{feasible} when \(P\ll\nu\), every \(f_k\) is \(P\)-integrable,
and \(E_P[f_k]=c_k\) for \(k=1,\ldots,m\). The \emph{feasible set} is
written \(\mathcal F_\nu(c)\). For feasible \(P\) with density
\(p=dP/d\nu\), the \emph{entropy relative to} \(\nu\) is
\[
 H_\nu(P)\defeq-\int_{\mathcal X}p\ln p\,d\nu ,
\]
whenever at least one of the positive and negative parts of \(p\ln p\)
has finite \(\nu\)-integral; it then lies in \([-\infty,+\infty]\).
The \emph{attainable moment set} is
\[
 \mathcal M_\nu\defeq\{E_P[f]: P\ll\nu,\ P(\mathcal X)=1,\
 f_k\in L^1(P)\text{ for all }k\}\subseteq\mathbb R^m ,
\]
so that \(\mathcal F_\nu(c)\neq\varnothing\) exactly when
\(c\in\mathcal M_\nu\). The \emph{value} of the problem is
\[
 v_\nu(c)\defeq\sup\{H_\nu(P):P\in\mathcal F_\nu(c),\ H_\nu(P)\text{ defined}\},
\]
with \(\sup\varnothing=-\infty\). A \emph{maximizer} is a feasible \(P\)
with defined entropy \(H_\nu(P)=v_\nu(c)\).
\index{maximum entropy!declared problem}\index{entropy!relative to a reference measure}
\end{hbdefinition}

\begin{hbconvention}{def:p15-w02b-p15-043-conventions}{Logarithms, null points and the source's notation}
Logarithms are natural and \(0\ln0=0\). A density is fixed up to
\(\nu\)-null sets, and the entropy does not depend on the version.
In the finite case, \(H_\nu(P)=-\sum_xP\{x\}\ln(P\{x\}/w(x))\); with
counting measure this is Jaynes's
\(H(p_1,\ldots,p_n)=-K\sum_ip_i\ln p_i\) with \(K=1\)
\cite[p.~622, eq.~(2-3); p.~623, after eq.~(2-13)]{Jaynes1957InformationTheory}.
On a countable space with counting measure every term \(-P\{x\}\ln P\{x\}\)
is nonnegative, so the entropy is always defined, with value in \([0,\infty]\).
Jaynes writes \(\langle f_r(x)\rangle\) for the expectation of \(f_r\) under the
assigned distribution, and \(\lambda_1,\ldots,\lambda_m\) for the multipliers;
both are kept below. Points of zero \(\nu\)-measure carry no feasible
probability and are removed in the finite case. A feasible distribution whose
entropy integral is undefined is not admitted as a competitor;
\cref{thm:p15-w02b-p15-043-gibbs} shows that this exclusion is empty
whenever the partition function is finite at some point.
\end{hbconvention}

The reference measure is part of the problem, not a neutral background.
On a finite space with counting measure it says that each listed point is
one possibility. On a continuum it fixes the coordinate in which ``spread
out'' is measured: Lebesgue measure in a length coordinate and Lebesgue
measure in an area coordinate are different reference measures for the same
set of objects, as \cref{thm:p15-w02b-p15-043-transport} makes exact.
Jaynes treats only the discrete case and states that the continuous case
is harder \cite[p.~622, n.~6]{Jaynes1957InformationTheory}. In his model of
closely spaced levels he replaces sums by integrals against a density of
states \(\rho(x)\), which is a reference measure in the present sense
\cite[p.~625]{Jaynes1957InformationTheory}.

\begin{hbdefinition}{def:p15-w02b-p15-043-partition}{Partition function, exponential family and relative entropy}
For \(\lambda\in\mathbb R^m\) write \(\lambda\cdot f=\sum_k\lambda_kf_k\) and
\[
 Z_\nu(\lambda)\defeq\int_{\mathcal X}e^{-\lambda\cdot f}\,d\nu\in(0,\infty],
 \qquad
 \Lambda_\nu\defeq\{\lambda\in\mathbb R^m:Z_\nu(\lambda)<\infty\}.
\]
\(Z_\nu\) is the \emph{partition function} and \(\Lambda_\nu\) its domain,
which may be empty. For \(\lambda\in\Lambda_\nu\) put
\(\lambda_0\defeq\ln Z_\nu(\lambda)\) and let \(Q_\lambda\) be the
probability measure with density
\[
 q_\lambda\defeq\frac{dQ_\lambda}{d\nu}
   =\exp\{-[\lambda_0+\lambda_1f_1+\cdots+\lambda_mf_m]\}.
\]
This is Jaynes's form (2-10) with (2-9) and (2-12), written relative to
\(\nu\) \cite[p.~623]{Jaynes1957InformationTheory}. The \emph{dual function}
for declared values \(c\) is \(\Psi_c(\lambda)\defeq\ln Z_\nu(\lambda)+\lambda\cdot c\)
on \(\Lambda_\nu\), extended by \(+\infty\) off \(\Lambda_\nu\).
The family is \emph{regular} when \(\Lambda_\nu\) is open and nonempty.
The statistics are \emph{minimal} for \(\nu\) when no
\(u\in\mathbb R^m\setminus\{0\}\) makes \(u\cdot f\) equal to a constant
\(\nu\)-almost everywhere.

For probability measures \(P,Q\ll\nu\) with densities \(p,q\), and
\(q>0\) \(\nu\)-almost everywhere, the \emph{relative entropy} is
\[
 D(P\Vert Q)\defeq\int_{\mathcal X}p\ln\frac pq\,d\nu .
\]
\index{partition function}\index{relative entropy}\index{maximum entropy!multipliers}
\end{hbdefinition}

\section{The optimizer and its qualifications}
\label{sec:p15-w02b-p15-043-optimizer}

\begin{hbtheorem}{thm:p15-w02b-p15-043-gibbs}{proved=annex:p15-w02b-p15-043-gibbs}{The Gibbs identity and the dual bound}
\begin{enumerate}
\item For probability measures \(P,Q\ll\nu\) with \(q>0\)
\(\nu\)-almost everywhere, \(D(P\Vert Q)\) is well defined in \([0,\infty]\),
and \(D(P\Vert Q)=0\) exactly when \(P=Q\).
\item Let \(\lambda\in\Lambda_\nu\) and \(P\in\mathcal F_\nu(c)\). Then
\(\ln q_\lambda\in L^1(P)\), \(H_\nu(P)\) is defined, and
\[
 H_\nu(P)=\lambda_0+\sum_{k=1}^m\lambda_kc_k-D(P\Vert Q_\lambda)
        =\Psi_c(\lambda)-D(P\Vert Q_\lambda).
\]
\item Consequently \(v_\nu(c)\leq\inf_{\lambda\in\Lambda_\nu}\Psi_c(\lambda)\),
whether or not any maximizer exists and whether or not
\(\mathcal F_\nu(c)\) is empty.
\end{enumerate}
\end{hbtheorem}

\begin{hbtheorem}{thm:p15-w02b-p15-043-optimality}{proved=annex:p15-w02b-p15-043-optimality}{A feasible exponential-family member is the unique maximizer}
Suppose \(\lambda\in\Lambda_\nu\), every \(f_k\) is \(Q_\lambda\)-integrable
and \(E_{Q_\lambda}[f_k]=c_k\) for all \(k\). Then \(Q_\lambda\) is the
unique maximizer of the declared problem and
\[
 v_\nu(c)=H_\nu(Q_\lambda)=\lambda_0+\lambda_1c_1+\cdots+\lambda_mc_m ,
\]
which is Jaynes's maximal entropy (2-13)
\cite[p.~623]{Jaynes1957InformationTheory}. Every other feasible \(P\)
satisfies \(H_\nu(P)=v_\nu(c)-D(P\Vert Q_\lambda)<v_\nu(c)\).
\end{hbtheorem}

\begin{hbcorollary}{thm:p15-w02b-p15-043-no-constraint}{proved=annex:p15-w02b-p15-043-no-constraint}{No constraints: normalized reference measure or unbounded entropy}
Let \(m=0\). If \(\nu(\mathcal X)<\infty\), the unique maximizer is
\(\nu/\nu(\mathcal X)\) and \(v_\nu=\ln\nu(\mathcal X)\). In the finite
case with weights \(w\) it gives probability \(w(x)/\sum_yw(y)\) to \(x\);
with counting measure on \(n\) points it is uniform with value \(\ln n\).
If \(\nu(\mathcal X)=\infty\), then \(v_\nu=+\infty\) and no feasible
distribution of finite entropy is a maximizer.
\end{hbcorollary}

The finite counting-measure case is Jaynes's observation that the principle
reduces to insufficient reason when only the enumeration of possibilities is
given \cite[p.~623]{Jaynes1957InformationTheory}; the weighted case shows
that a declared reference measure, not ignorance alone, fixes the answer.

\begin{hbtheorem}{thm:p15-w02b-p15-043-uniqueness}{proved=annex:p15-w02b-p15-043-uniqueness}{Strict concavity and uniqueness}
The feasible set \(\mathcal F_\nu(c)\) is convex. If
\(P_1,P_2\in\mathcal F_\nu(c)\) have finite entropies, then
\(M=\tfrac12(P_1+P_2)\in\mathcal F_\nu(c)\) has defined entropy and
\[
 H_\nu(M)\geq\tfrac12\bigl(H_\nu(P_1)+H_\nu(P_2)\bigr),
\]
with equality only if \(P_1=P_2\). Hence a declared problem with finite
value has at most one maximizer. When \(v_\nu(c)=+\infty\) no uniqueness
is asserted.
\end{hbtheorem}

\begin{hblemma}{thm:p15-w02b-p15-043-partition}{proved=annex:p15-w02b-p15-043-partition}{Properties of the partition function}
\begin{enumerate}
\item \(\Lambda_\nu\) is convex, \(\ln Z_\nu\) is convex on it, and
\(Z_\nu\) is lower semicontinuous on \(\mathbb R^m\) as a
\([0,\infty]\)-valued function.
\item Let \(\lambda\) be an interior point of \(\Lambda_\nu\). Every
product \(f_jf_k\) and every \(f_k\) is \(Q_\lambda\)-integrable. For each
\(u\in\mathbb R^m\) the function \(t\mapsto\ln Z_\nu(\lambda+tu)\) is
twice differentiable at \(t=0\) with first derivative
\(-E_{Q_\lambda}[u\cdot f]\) and second derivative
\(\operatorname{Var}_{Q_\lambda}(u\cdot f)\). In particular
\[
 \langle f_r\rangle=-\frac{\partial}{\partial\lambda_r}\ln Z_\nu,
 \qquad
 \langle f_r^2\rangle-\langle f_r\rangle^2=\frac{\partial^2}{\partial\lambda_r^2}\ln Z_\nu ,
\]
Jaynes's (2-11) and (2-14) \cite[p.~623]{Jaynes1957InformationTheory}.
\item If \(f\) is minimal, these second derivatives are positive for
\(u\neq0\), and \(\Psi_c\) is strictly convex on every segment inside the
interior of \(\Lambda_\nu\).
\item If \(f\) is not minimal and \(u\neq0\) makes \(u\cdot f\) constant
\(\nu\)-almost everywhere, then \(Q_{\lambda+su}=Q_\lambda\) whenever
\(\lambda,\lambda+su\in\Lambda_\nu\).
\end{enumerate}
\end{hblemma}

\begin{hbtheorem}{thm:p15-w02b-p15-043-interior}{proved=annex:p15-w02b-p15-043-interior}{Moments of positive densities are interior}
Let \(f\) be minimal. If \(P\in\mathcal F_\nu(c)\) has a density that is
positive \(\nu\)-almost everywhere and
\(E_P[f_k^2]<\infty\) for all \(k\), then \(c\) is an interior point of
\(\mathcal M_\nu\).
\end{hbtheorem}

\begin{hbtheorem}{thm:p15-w02b-p15-043-existence}{proved=annex:p15-w02b-p15-043-existence}{Existence and uniqueness of multipliers}
Let the family be regular and \(f\) minimal. Then the mean map
\(\lambda\mapsto E_{Q_\lambda}[f]\) is a bijection from \(\Lambda_\nu\)
onto the interior of \(\mathcal M_\nu\). Consequently:
\begin{enumerate}
\item if \(c\) is interior to \(\mathcal M_\nu\), the multiplier equations
\(-\partial\ln Z_\nu/\partial\lambda_k=c_k\) \((k=1,\ldots,m)\) have exactly
one solution \(\lambda(c)\in\Lambda_\nu\); \(Q_{\lambda(c)}\) is the unique
maximizer, and
\(v_\nu(c)=\Psi_c(\lambda(c))=\min_{\Lambda_\nu}\Psi_c\);
\item if \(c\in\mathcal M_\nu\) is not interior, the multiplier equations
have no solution in \(\Lambda_\nu\), and no maximizer has the form
\(Q_\lambda\);
\item if \(c\notin\mathcal M_\nu\), the feasible set is empty and the
multiplier equations have no solution.
\end{enumerate}
\end{hbtheorem}

\begin{hbtheorem}{thm:p15-w02b-p15-043-finite}{proved=annex:p15-w02b-p15-043-finite}{The finite case, including boundary values}
In the finite case:
\begin{enumerate}
\item \(\mathcal M_\nu\) is the convex hull of \(\{f(x):x\in\mathcal X\}\),
and \(\Lambda_\nu=\mathbb R^m\).
\item For every \(c\in\mathcal M_\nu\) a maximizer exists, it is unique and
\(v_\nu(c)\leq\ln\sum_xw(x)\).
\item The support of the maximizer is the set
\(S(c)\) of points that receive positive probability from at least one
feasible distribution.
\item If \(f\) is minimal, the following are equivalent: \(c\) is interior to
\(\mathcal M_\nu\); the multiplier equations have a solution; the
maximizer gives positive probability to every point of \(\mathcal X\).
\end{enumerate}
\end{hbtheorem}

Item 3 is the exact content of Jaynes's remark that the maximum-entropy
distribution assigns positive weight to every situation the given information
does not absolutely exclude \cite[p.~623]{Jaynes1957InformationTheory}.
The excluded situations are those to which the constraints themselves assign
zero probability in every feasible distribution. At a boundary value of
\(c\) such points exist, and the form (2-10) with finite multipliers is then
unavailable.

\begin{hbtheorem}{thm:p15-w02b-p15-043-failures}{proved=annex:p15-w02b-p15-043-failures}{Nonexistence and nonattainment}
\begin{enumerate}
\item \emph{Boundary value.} On \(\mathcal X=\{0,1,2\}\) with counting
measure, \(f(x)=x\) and \(c=2\), the feasible set is \(\{\delta_2\}\),
the maximizer is \(\delta_2\) with entropy \(0\), and no
\(\lambda\in\mathbb R\) solves the multiplier equation.
\item \emph{Unbounded support with too few constraints.} On \(\mathbb R\)
with Lebesgue measure, \(f(x)=x\) and any \(c\), the partition function is
infinite everywhere, \(\Lambda_\nu=\varnothing\), and \(v_\nu(c)=+\infty\);
no feasible distribution of finite entropy is a maximizer.
\item \emph{Partition function infinite on part of the line.} On
\([0,\infty)\) with \(\nu(dx)=(1+x)^{-3}\,dx\) and \(f(x)=x\):
\(\Lambda_\nu=[0,\infty)\), which is not open;
\(\mathcal M_\nu=(0,\infty)\); and every \(Q_\lambda\) has mean at most
\(1\). For \(c>1\), \(c\) is interior to \(\mathcal M_\nu\), the
multiplier equation has no solution, \(v_\nu(c)=-\ln2=\Psi_c(0)
=\inf_{\Lambda_\nu}\Psi_c\), and the value is not attained.
\item \emph{Unattainable value.} In the setting of item 1 with \(c=3\), the
feasible set is empty, \(v_\nu(3)=-\infty\), and
\(\Psi_3(\lambda)\to-\infty\) as \(\lambda\to-\infty\).
\end{enumerate}
\end{hbtheorem}

Item 3 shows why regularity is a hypothesis of
\cref{thm:p15-w02b-p15-043-existence} rather than a technicality: the
dual bound is still exact, but its minimizer sits on the edge of
\(\Lambda_\nu\), where the corresponding \(Q_0\) violates the constraint.

\section{Multipliers and normalization}
\label{sec:p15-w02b-p15-043-multipliers}

Jaynes introduces one multiplier \(\lambda_r\) for each given average and a
further constant \(\lambda_0\) for normalization, obtains (2-10), and fixes
the constants by (2-11) and (2-12) \cite[p.~623]{Jaynes1957InformationTheory}.
The results above separate what each of these steps needs.

\paragraph{Normalization.}
For every \(\lambda\in\Lambda_\nu\), \(\lambda_0=\ln Z_\nu(\lambda)\) makes
\(q_\lambda\) a probability density. This uses no constraint value and no
optimizer. Outside \(\Lambda_\nu\) the exponential weight cannot be
normalized, and there is no member \(Q_\lambda\) at all.

\paragraph{The multiplier equations.}
The equations \(-\partial\ln Z_\nu/\partial\lambda_k=c_k\) are meaningful at
interior points of \(\Lambda_\nu\), where they say
\(E_{Q_\lambda}[f_k]=c_k\) (\cref{thm:p15-w02b-p15-043-partition}).
A solution, if one exists, produces a feasible \(Q_\lambda\) and therefore,
by \cref{thm:p15-w02b-p15-043-optimality}, the maximizer. Whether a solution
exists is a separate question, answered by
\cref{thm:p15-w02b-p15-043-existence,thm:p15-w02b-p15-043-finite,thm:p15-w02b-p15-043-failures}:
under regularity and minimality there is exactly one solution when \(c\) is
interior to \(\mathcal M_\nu\), and none otherwise.

\paragraph{Multipliers as dual minimizers.}
For every \(\lambda\in\Lambda_\nu\) the number \(\Psi_c(\lambda)\) bounds the
entropy of every feasible distribution
(\cref{thm:p15-w02b-p15-043-gibbs}). This bound does not presuppose that the
feasible set is nonempty or that the supremum is attained. Solving the
multiplier equations is the same as finding a stationary point of the convex
function \(\Psi_c\). Under regularity and minimality the stationary point is
the minimizer and the bound is attained. Without regularity the infimum of
the bound may still equal the value but be approached only at the edge of
\(\Lambda_\nu\), as in \cref{thm:p15-w02b-p15-043-failures}.

\paragraph{When the feasible set is empty.}
If \(c\notin\mathcal M_\nu\), then \(v_\nu(c)=-\infty\) and there is nothing
to maximize. The dual bound still holds but conveys nothing about a
distribution. The function \(\Psi_c\) may be unbounded below, as in item 4 of
\cref{thm:p15-w02b-p15-043-failures}. Declared values
outside the attainable set signal that the declarations are inconsistent
with the reference measure; they do not produce a degenerate optimizer.

\paragraph{Redundant statistics.}
If the statistics are not minimal, the multipliers are not unique but the
distribution is: shifting \(\lambda\) along a direction \(u\) with \(u\cdot f\)
constant leaves \(Q_\lambda\) unchanged
(\cref{thm:p15-w02b-p15-043-partition}). The multipliers then have no
individual interpretation; only the distribution does.

\begin{hbmethod}{def:p15-w02b-p15-043-method}{Maximum-entropy assignment under declared constraints}
\begin{enumerate}
\item Declare the space \(\mathcal X\), the reference measure \(\nu\)
(with the coordinate in which it is uniform), the statistics
\(f_1,\ldots,f_m\) and the values \(c_1,\ldots,c_m\).
\item Check feasibility: exhibit one feasible distribution, or show
\(c\notin\mathcal M_\nu\) and stop with an empty problem.
\item Compute \(Z_\nu\) and its domain \(\Lambda_\nu\). Note whether the
family is regular and whether the statistics are minimal; remove redundant
statistics if not.
\item Solve \(-\partial\ln Z_\nu/\partial\lambda_k=c_k\) in the interior of
\(\Lambda_\nu\). If a solution exists, check \(E_{Q_\lambda}[f_k]=c_k\)
directly; \(Q_\lambda\) is then the unique maximizer, with value
\(\lambda_0+\sum_k\lambda_kc_k\).
\item If no solution exists, classify the failure: boundary value
(finite case: the maximizer exists on the support \(S(c)\)); empty
domain or unbounded value; or a finite value approached but not attained.
Report the classification instead of a distribution.
\item State the result as the solution of the declared problem. Any
inferential use is conditional on the premises in
\cref{sec:p15-w02b-p15-043-jaynes}, including the choice of \(\nu\).
\end{enumerate}
\end{hbmethod}

\begin{hbtheorem}{thm:p15-w02b-p15-043-transport}{proved=annex:p15-w02b-p15-043-transport}{Reparametrization and rescaling of the reference measure}
\begin{enumerate}
\item Let \(T:(\mathcal X,\mathcal A)\to(\mathcal Y,\mathcal B)\) be a
bijection with \(T\) and \(T^{-1}\) measurable, and let \(T_*\nu\) be the
image measure. Then \(P\mapsto T_*P\) maps \(\mathcal F_\nu(c)\) for the
statistics \(f\) bijectively onto \(\mathcal F_{T_*\nu}(c)\) for the
statistics \(f\circ T^{-1}\), and \(H_{T_*\nu}(T_*P)=H_\nu(P)\). Maximizers
correspond under \(T\).
\item For a constant \(a>0\), \(H_{a\nu}(P)=H_\nu(P)+\ln a\) for every
\(P\ll\nu\) with defined entropy; the maximizers for \(\nu\) and
\(a\nu\) coincide.
\item On \((0,\infty)\) with \(T(x)=x^2/2\), the image of \(x\,dx\) is
Lebesgue measure \(du\) on \((0,\infty)\), while the image of \(dx\) is
\((2u)^{-1/2}\,du\).
\end{enumerate}
\end{hbtheorem}

Part 1 says the problem is coordinate-free once the reference measure is
carried along. Part 3 shows what changes when it is not carried along:
declaring Lebesgue measure in \(x\) and declaring Lebesgue measure in
\(u=x^2/2\) are two different reference measures on the same set. This is the
precise form of the representation dependence in the factory example of
\cref{sec:p15-rationality-4-1}: uniformity in side length and uniformity in
face area are maximum-entropy assignments for different reference measures.
Part 2 is why Jaynes's constant \(1/\epsilon\) in \(\rho(x)=x^n/\epsilon\) does not
affect his estimates \cite[p.~625]{Jaynes1957InformationTheory}.

\section{Jaynes's inferential argument}
\label{sec:p15-w02b-p15-043-jaynes}

The paper's normative claim is that, in making inferences from partial
information, one should use the distribution of maximum entropy subject to
what is known, because any other choice assumes information one does not
have \cite[p.~623]{Jaynes1957InformationTheory}. Reconstructed, the argument
has the following premises.

\begin{enumerate}
\item[(J1)] \emph{Underdetermination.} The variable takes discrete values
\(x_1,\ldots,x_n\); all that is known is the expectation of a function
\(f(x)\), eq.~(2-1), together with normalization, eq.~(2-2). These two
conditions leave \(n-2\) further conditions undetermined, so the probabilities,
and the expectation of another function \(g(x)\), are not fixed by the data
\cite[p.~622]{Jaynes1957InformationTheory}. In statistical mechanics the same
situation arises: macroscopic data are far too few to determine probabilities
of individual quantum states \cite[p.~622, n.~5]{Jaynes1957InformationTheory}.
\item[(J2)] \emph{Critique of insufficient reason.} Laplace's rule assigns
equal probabilities when there is no reason to think otherwise. Jaynes grants
it where an evident symmetry makes the events equally possible, but holds it
otherwise as arbitrary as any other assumption; for continuous quantities it
generates paradoxes because ``equally possible'' changes under a change of
variables. What was missing was a constructive criterion for preferring one
distribution among those that agree equally well with the information
\cite[p.~622]{Jaynes1957InformationTheory}.
\item[(J3)] \emph{Interpretation.} Jaynes distinguishes an objective school,
for which a probability is a frequency property testable by fluctuations,
from a subjective school, for which it expresses a state of knowledge and is
tested by whether it represents that knowledge. A problem of type (J1) belongs
to the second kind, so the subjective view is adopted for it. He also holds
that physics needs both views and calls the subjective one the broader
\cite[p.~622]{Jaynes1957InformationTheory}.
\item[(J4)] \emph{A unique measure of uncertainty.} Following Shannon, the
amount of uncertainty represented by a discrete distribution is
\(H=-K\sum_ip_i\ln p_i\) with \(K>0\), uniquely fixed by continuity,
monotonicity in the number of equally likely alternatives and a composition
law \cite[p.~622, eq.~(2-3); p.~630, eqs.~(A-1)--(A-5)]{Jaynes1957InformationTheory}.
Jaynes sketches the derivation and refers to Shannon for the step that
\(A(n)=K\ln n\) is the only solution of (A-3);
\cref{thm:p15-w02b-p15-043-shannon} supplies that step and the rest of the
derivation. He prefers this measure to \(-\sum p_i^2\) because it alone
satisfies the composition law, and because stationary points of the latter
need not satisfy \(p_i\geq0\) \cite[p.~623]{Jaynes1957InformationTheory}.
\item[(J5)] \emph{Least bias.} An assignment should agree with the given
information and avoid bias; choosing the distribution of maximal uncertainty
does this, since a smaller entropy would represent information not possessed.
The maximum-entropy distribution is then defended positively, as maximally
noncommittal about missing information, not negatively, as the absence of a
reason to think otherwise \cite[pp.~622--623]{Jaynes1957InformationTheory}.
\end{enumerate}

\begin{hbtheorem}{thm:p15-w02b-p15-043-shannon}{proved=annex:p15-w02b-p15-043-shannon}{The uncertainty measure from Jaynes's Appendix~A conditions}
For each \(n\geq1\) let \(H\) be a real function on the probability vectors
\((p_1,\ldots,p_n)\). Suppose:
(1) \(H\) is continuous on each such simplex;
(2) \(A(n)\defeq H(1/n,\ldots,1/n)\) is strictly increasing in \(n\);
(3) for every grouping of \(1,\ldots,n\) into consecutive groups
\(G_1,\ldots,G_r\) with group probabilities
\(w_j=\sum_{i\in G_j}p_i>0\),
\[
 H(p_1,\ldots,p_n)=H(w_1,\ldots,w_r)+\sum_{j=1}^rw_j\,H\bigl((p_i/w_j)_{i\in G_j}\bigr).
\]
Then there is a constant \(K>0\) with
\(H(p_1,\ldots,p_n)=-K\sum_ip_i\ln p_i\) for all \(n\) and all probability
vectors.
\end{hbtheorem}

Condition (3) is the composition law (A-1)
\cite[p.~630]{Jaynes1957InformationTheory}. Under (J1)--(J5), with the
expectation information as constraints and the enumeration of states as the
reference measure, the conclusion follows: \emph{the distribution to be used
is the solution of the declared problem}, namely (2-10) when the multiplier
equations are solvable. The mathematics of the previous sections adds
qualifications that the paper does not state: the given averages must be
attainable; at boundary values the solution exists but is not of the form
(2-10); and outside the finite case a solution may fail to exist. The
conclusion is conditional at each point. If (J3) is replaced by a frequency
interpretation, the optimizer carries no claim about long-run frequencies.
If the composition law in (J4) is rejected, the uniqueness of \(H\) lapses.
If the reference measure is changed, the conclusion changes with it
(\cref{thm:p15-w02b-p15-043-transport,def:p15-w02b-p15-043-reference-example}).

\paragraph{Application to statistical mechanics and its stated limits.}
With energy levels \(E_i(\alpha_1,\alpha_2,\ldots)\) and only \(\langle E\rangle\)
given, the special case of (2-10) is the Boltzmann distribution. Identifying
\(\lambda_1=1/kT\) and the free energy recovers the usual relations (3-1)--(3-4).
If particle numbers are also uncertain and their averages are given, the
grand canonical ensemble results, (3-5)--(3-8) \cite[pp.~623--624]{Jaynes1957InformationTheory}.
Jaynes bounds this application in several explicit ways.
\begin{itemize}
\item It concerns prediction of equilibrium properties, using probabilities of
stationary states, from data on quantities observed to be constant in time;
density-matrix refinements and irreversible processes are deferred to later
work \cite[p.~621; p.~624]{Jaynes1957InformationTheory}. Only the discrete
case is treated \cite[p.~622, n.~6]{Jaynes1957InformationTheory}.
\item The only point of contact with physics is the enumeration of the
mutually exclusive possible states; a new advance in knowledge changes the
inference equations only by changing that enumeration
\cite[p.~625]{Jaynes1957InformationTheory}. In the terms of this chapter, the
physics enters through \(\mathcal X\) and \(\nu\).
\item The information alone is far too small for reliable predictions. Definite
predictions arise only when the distributions of the measured macroscopic
quantities are sharply peaked, as they are for very many degrees of freedom; the
model (3-9)--(3-14) shows relative widths falling like \(n^{-1/2}\)
\cite[pp.~625--626]{Jaynes1957InformationTheory}. Jaynes calls the requirement
that almost all appreciably weighted states have the same macroscopic behavior
the principle of macroscopic uniformity, and locates the objective content of
the calculations there rather than in the probabilities
\cite[p.~627]{Jaynes1957InformationTheory}.
\item A sharp prediction contradicted by experiment is read as evidence that the
enumeration of possible states was incorrect, not that the information was
insufficient \cite[p.~627]{Jaynes1957InformationTheory}.
\item Interpretation, inferring microscopic causes from observed behavior, needs
probabilities in the frequency sense and a different theory, which Jaynes calls
objective statistical mechanics \cite[p.~627]{Jaynes1957InformationTheory}.
\item The generalization in which temperature-like quantities are attached to any
conserved, exchangeable quantity rests on a heat-bath factorization argument that
Jaynes says lacks rigor without the density matrix
\cite[pp.~627--628, n.~18]{Jaynes1957InformationTheory}.
\end{itemize}
The paper concludes that, in prediction, maximizing entropy is a method of
reasoning, not a law of physics \cite[p.~630]{Jaynes1957InformationTheory}.

\begin{hbwarning}{def:p15-w02b-p15-043-not-established}{What the optimization does not establish}
A maximizer is the solution of a declared problem. It is not shown to be
the true distribution, and it acquires no evidential support from being the
maximizer. Its use as an inference rests on (J3)--(J5) and on the declared
reference measure and constraints. In Jaynes's own account, success in
statistical mechanics comes from sharp peaks and correct enumeration of states,
not from the entropy criterion alone \cite[pp.~625--627]{Jaynes1957InformationTheory}.
Entropies relative to different reference measures are not comparable numbers:
rescaling \(\nu\) by \(a\) shifts every entropy by \(\ln a\)
(\cref{thm:p15-w02b-p15-043-transport}).
\end{hbwarning}

\section{Worked examples: changing the reference measure and the constraint}
\label{sec:p15-w02b-p15-043-examples}

The following family is Jaynes's density-of-states model with the constant
\(1/\epsilon\) removed (which \cref{thm:p15-w02b-p15-043-transport} permits);
the exponent \(n\) is allowed to be any real number above \(-1\).

\begin{hbtheorem}{thm:p15-w02b-p15-043-power-family}{proved=annex:p15-w02b-p15-043-power-family}{Maximum entropy against the reference measures \(x^n\,dx\)}
Let \(\mathcal X=(0,\infty)\), \(n>-1\) and \(\nu_n(dx)=x^n\,dx\).
\begin{enumerate}
\item With \(f(x)=x\) and \(c>0\): \(\Lambda_{\nu_n}=(0,\infty)\),
\(Z(\lambda)=\Gamma(n+1)\lambda^{-(n+1)}\), the unique multiplier is
\(\lambda=(n+1)/c\), and the unique maximizer has Lebesgue density
\(\lambda^{n+1}x^ne^{-\lambda x}/\Gamma(n+1)\), with value
\(\ln\Gamma(n+1)-(n+1)\ln\lambda+n+1\) and
\(E[x^2]=(n+2)c^2/(n+1)\).
\item With \(f(x)=x^2\) and \(s>0\): \(\Lambda_{\nu_n}=(0,\infty)\),
\(Z(\lambda)=\Gamma(\tfrac{n+1}2)/(2\lambda^{(n+1)/2})\), the unique
multiplier is \(\lambda=(n+1)/(2s)\), and the unique maximizer has
Lebesgue density \(2\lambda^{(n+1)/2}x^ne^{-\lambda x^2}/\Gamma(\tfrac{n+1}2)\),
with value \(\ln Z(\lambda)+\lambda s\) and
\[
 E[x]=\frac{\Gamma(\tfrac{n+2}2)}{\Gamma(\tfrac{n+1}2)}
      \Bigl(\frac{2s}{n+1}\Bigr)^{1/2}.
\]
\end{enumerate}
For integer \(n\geq0\) these are Jaynes's problems (A) and (B), and the two
expectation formulas are his (3-10) and (3-11)
\cite[p.~625]{Jaynes1957InformationTheory}.
\end{hbtheorem}

\begin{hbexample}{def:p15-w02b-p15-043-reference-example}{Same constraint, two reference measures}
Declare \(x>0\) and \(E[x]=2\).
\begin{enumerate}
\item[(i)] Reference \(\nu_0=dx\), Lebesgue measure in \(x\). By
\cref{thm:p15-w02b-p15-043-power-family} with \(n=0\), \(c=2\):
\(Z(\lambda)=1/\lambda\), \(\lambda=1/2\), and the maximizer is the exponential
density \(\tfrac12e^{-x/2}\). Check: \(\int_0^\infty x\cdot\tfrac12e^{-x/2}dx=2\);
\(Z(1/2)=2<\infty\); so \cref{thm:p15-w02b-p15-043-optimality} applies and the
value is \(\ln2+1\approx1.6931\). It gives \(E[x^2]=8\) and
\(P(x\leq1)=1-e^{-1/2}\approx0.3935\).
\item[(ii)] Reference \(\nu_1=x\,dx\), which is Lebesgue measure in the coordinate
\(u=x^2/2\) (\cref{thm:p15-w02b-p15-043-transport}). With \(n=1\), \(c=2\):
\(Z(\lambda)=\lambda^{-2}\), \(\lambda=1\), and the maximizer has Lebesgue density
\(xe^{-x}\) (density \(e^{-x}\) relative to \(\nu_1\)). Check:
\(\int_0^\infty x\cdot xe^{-x}dx=2\); the value is \(0-0+2=2\). It gives
\(E[x^2]=6\) and \(P(x\leq1)=1-2e^{-1}\approx0.2642\). In the coordinate \(u\)
the same distribution has Lebesgue density \(e^{-\sqrt{2u}}\), the
maximum-entropy density for Lebesgue measure in \(u\) under
\(E[\sqrt{2u}]=2\).
\end{enumerate}
The declared information is identical in (i) and (ii); the predictions for
\(x^2\) and for the event \(x\leq1\) differ. The two values \(1.6931\) and
\(2\) are entropies relative to different measures and are not compared. The
difference between the answers is produced entirely by the declared reference
measure, that is, by the coordinate in which one declares indifference; nothing
in the declared information favors either.
\end{hbexample}

\begin{hbexample}{def:p15-w02b-p15-043-constraint-example}{Same reference measure, a changed constraint}
Keep \(\nu_1=x\,dx\) and replace the declaration \(E[x]=2\) by
\(E[x^2]=6\), the value that (ii) predicts.
\begin{enumerate}
\item[(iii)] By \cref{thm:p15-w02b-p15-043-power-family} with \(n=1\), \(s=6\):
\(Z(\lambda)=1/(2\lambda)\), \(\lambda=1/6\), and the maximizer has Lebesgue
density \(\tfrac13xe^{-x^2/6}\). Check: the substitution \(y=x^2/6\) gives
\(\int_0^\infty x^2\cdot\tfrac13xe^{-x^2/6}dx=6\int_0^\infty ye^{-y}dy=6\), and
\(Z(1/6)=3<\infty\), so it is the unique maximizer, with value
\(\ln3+1\approx2.0986\).
It gives \(E[x]=\Gamma(\tfrac32)\,6^{1/2}/\Gamma(1)=\tfrac12\sqrt{6\pi}\approx2.1708\)
and \(P(x\leq1)=1-e^{-1/6}\approx0.1535\).
\end{enumerate}
The distribution \(xe^{-x}\) of (ii) is feasible for (iii), since its second
moment is \(6\). By \cref{thm:p15-w02b-p15-043-gibbs} its entropy relative to
\(\nu_1\) is the value of (iii) minus \(D\); the entropy is \(2\), so
\(D=\ln3+1-2=\ln3-1\approx0.0986>0\), as uniqueness requires.
The optimizer of (iii) is not feasible for (ii), since its mean is
\(\tfrac12\sqrt{6\pi}\neq2\). Starting from \(E[x]=2\), predicting \(E[x^2]=6\), and
then treating that prediction as the declared information yields a different
distribution and a different estimate of \(x\). This is the asymmetry of
Jaynes's two regression lines \cite[pp.~625--626, Fig.~1]{Jaynes1957InformationTheory}.
The second problem uses no new evidence; it uses a different declaration,
and its answer is a property of that declaration.
\end{hbexample}

\begin{hbexample}{def:p15-w02b-p15-043-copper-example}{The asymmetric copper coefficient, computed}
The guide's interval \([16.40,16.92]\) with declared mean \(16.52\), in units
\(10^{-6}\,{}^\circ\mathrm C^{-1}\), is the problem of
\cref{thm:p15-measurement-asymmetric-model} with Lebesgue reference measure in
the coefficient itself, \(a=b_-=0.12\) and \(b=b_+=0.40\). In the shifted
coordinate \(T\in[-a,b]\) one has
\(Z(\lambda)=(e^{\lambda a}-e^{-\lambda b})/\lambda\) for \(\lambda\neq0\), and
integration by parts gives, for \(\lambda\neq0\),
\[
 E[T]=0\iff \lambda\bigl(be^{-\lambda b}+ae^{\lambda a}\bigr)=e^{\lambda a}-e^{-\lambda b}.
\]
At \(\lambda=0\) the density is uniform with \(E[T]=(b-a)/2\neq0\), so the
multiplier is not zero. The left side minus the right side is about
\(-0.0040\) at \(\lambda=7.7\) and \(+0.0187\) at \(\lambda=7.8\), and the
multiplier is unique by \cref{thm:p15-measurement-asymmetric-model}. It is \(\lambda\approx7.7177\) per unit of
the coefficient; the declared mean below the midpoint puts the larger density at
the lower end. The density is \(e^{-\lambda(t-16.52)}/Z(\lambda)\) on the
interval, with \(Z\approx0.3212\), so its values at the two ends are about
\(7.860\) and \(0.142\). The variance formula of
\cref{thm:p15-measurement-asymmetric-model} gives
\(0.048-0.28/7.7177\approx0.01172\), a standard deviation of about \(0.108\),
and probability about \(0.615\) below the declared mean.
Without the mean declaration the maximizer is uniform, with mean \(16.66\),
standard deviation \(0.52/\sqrt{12}\approx0.150\) and probability
\(0.12/0.52\approx0.231\) below \(16.52\). Adding the declaration lowers the
maximal entropy from \(\ln0.52\approx-0.654\) to
\(\ln Z\approx-1.136\), as adding a constraint must. Both are results of
declared problems with the same reference measure. Neither computation adds
evidence about the copper sample; the difference records the added declared
mean.
\end{hbexample}

\section{Limitations and failure cases}
\label{sec:p15-w02b-p15-043-limits}

The method fails, or must report something other than a distribution, in the
cases proved in \cref{thm:p15-w02b-p15-043-failures,thm:p15-w02b-p15-043-no-constraint}:
values outside the attainable set; boundary values, where the optimizer
exists in the finite case but has no multipliers; an unbounded value when the
reference measure is infinite and the constraints do not confine it; and a
finite value that is not attained when the partition-function domain is not
open. A procedure that always returns a formula of the form (2-10) hides
these cases.

The entropy of a continuous distribution has no meaning apart from its
reference measure; \cref{def:p15-w02b-p15-043-reference-example} shows two
different answers from identical expectation information. The declaration of
\(\nu\) therefore carries the same evidential burden as any other premise.
Jaynes himself restricts the paper to the discrete case
\cite[p.~622, n.~6]{Jaynes1957InformationTheory}.

Expectation constraints are a special form of information. They say that an
average equals a number, not that a sample average was observed, and the
method has no rule for converting a finite set of observations into exact
expectation constraints. When the information is not of this form, the
declared problem is a model of the information, and its adequacy is a
separate question.

\section{Relationships}
\label{sec:p15-w02b-p15-043-relations}

\begin{itemize}
\item \cref{thm:p15-rat41-indifference} is the case \(m=0\) of
\cref{thm:p15-w02b-p15-043-no-constraint} with counting measure; the factory
example in \cref{sec:p15-rationality-4-1} is the reference-measure dependence of
\cref{thm:p15-w02b-p15-043-transport}.
\item \cref{thm:p15-measurement-asymmetric-model} is the case of one mean on a
bounded interval with Lebesgue reference measure; there \(\Lambda_\nu=\mathbb R\)
and \cref{thm:p15-w02b-p15-043-existence} applies to every interior mean.
\item \cref{thm:p15-rmdl-gibbs} is the countable form of the inequality
\(D\geq0\) in \cref{thm:p15-w02b-p15-043-gibbs}.
\item Shannon's own statement of the characterization used in (J4) is
\cref{thm:p15-w02b-p15-040-characterization};
\cref{thm:p15-w02b-p15-043-shannon} is the form of Jaynes's Appendix~A.
\item Entropy relative to a reference measure, and its behavior under a change
of reference measure or of coordinates, is developed in
\cref{sec:p15-w02b-p15-042-continuous-entropy}. A criterion for attainment from
the dual function alone, the effect on the maximizer of a reference measure with
a nonconstant density, and the factory example with a mean constraint are in
\cref{sec:p15-w02b-p15-044-maxent-existence}.
\item Reparametrization in \cref{thm:p15-w02b-p15-043-transport} uses image
measures as in \cref{thm:p15-pushforward-laws}; the reparameterization contrast in
\cref{thm:p15-bayes-transformation} is the same distinction for priors.
\end{itemize}

\paragraph{Source.}
E.~T.~Jaynes, ``Information Theory and Statistical Mechanics,'' \emph{Physical
Review} 106 (1957), 620--630: Sections~1--6 and Appendix~A
\cite{Jaynes1957InformationTheory}. The general reference-measure results,
existence conditions and failure cases are not in that paper; they are stated and
proved here.
